\documentclass{article}

\usepackage{geometry}
\geometry{a4paper}
%\usepackage[UTF8, heading = false, scheme = plain]{ctex}%格式
\usepackage{ctex}
\usepackage{graphicx} %添加图片
\usepackage{amsthm}
\usepackage{appendix}
\usepackage{amsmath}
\usepackage{tocloft}
\usepackage{listings}
\renewcommand{\vec}[1]{\boldsymbol{#1}} % 生产粗体向量，而不是带箭头的向量
\usepackage{amssymb}
\usepackage{booktabs} % excel导出的大表格
\usepackage{fontspec}
\usepackage{xcolor}

 \usepackage{caption}
 \usepackage{float}
 \usepackage{wrapfig}
\usepackage{mathrsfs}
\usepackage{epstopdf}

\captionsetup[figure]{labelfont={bf},name={Fig.},labelsep=period}
 
% \setmainfont{Times New Roman}

\newcommand{\xiaoer}{\fontsize{18pt}{18pt}\selectfont}
\newcommand{\yihao}{\fontsize{26pt}{26pt}\selectfont}
\newcommand{\xiaosan}{\fontsize{15pt}{22pt}\selectfont}
% \renewcommand{\baselinestretch}{1.8}    %两倍行距
\renewcommand{\contentsname}{\hfill\bfseries\Large Contents\hfill}   
\renewcommand{\cftaftertoctitle}{\hfill}
% \definecolor{grey}{rgb}{0.8,0.8,0.8}
% \definecolor{darkgreen}{rgb}{0,0.3,0}
% \definecolor{darkblue}{rgb}{0,0,0.3}
% \def\lstbasicfont{\fontfamily{pcr}\selectfont\footnotesize}
% \lstset{%
%    % numbers=left,
%    % numberstyle=\small,%
%    	language=Matlab,
%     showstringspaces=false,
%     showspaces=false,%
%     breaklines=true,%
%     tabsize=4,%
%     frame=lines,%
%     basicstyle={\footnotesize\lstbasicfont},%
%     keywordstyle=\color{darkblue}\bfseries,%
%     identifierstyle=,%
%     commentstyle=\color{darkgreen},%\itshape,%
%     stringstyle=\color{black},%
% }

\usepackage{color} %red, green, blue, yellow, cyan, magenta, black, white
\definecolor{mygreen}{RGB}{28,172,0} % color values Red, Green, Blue
\definecolor{mylilas}{RGB}{170,55,241}
% \def\lstbasicfont{\fontfamily{pcr}\selectfont\footnotesize}
\lstset{language=Matlab,%
    %basicstyle=\color{red},
    breaklines=true,%
    tabsize=4,%
    morekeywords={matlab2tikz},
    morekeywords=[2]{1}, keywordstyle=[2]{\color{black}},
    showstringspaces=false,%without this there will be a symbol in the places where there is a space
    numbers=left,%
    numberstyle={\tiny \color{black}},% size of the numbers
    numbersep=9pt, % this defines how far the numbers are from the text
    emph=[1]{for,end,break},emphstyle=[1]\color{red}, %some words to emphasise
    %emph=[2]{word1,word2}, emphstyle=[2]{style},   
    frame=lines,% 
    % basicstyle={\footnotesize\lstbasicfont},%
    keywordstyle=\color{blue},%
    identifierstyle=\color{black},%
    stringstyle=\color{mylilas},
    commentstyle=\color{mygreen},%
}




\begin{document}

\begin{figure}[htbp]
\flushright
\includegraphics[width = 0.25 \linewidth]{./image/log.png}
\end{figure}

\vspace{40pt}

\centerline{\xiaoer\textbf{Laboratory Report of Digital Signal Processing}}

\vspace{10pt}
\centerline{\xiaosan Lab1.XXXX} % Times New Roman，小二


\vspace{80pt}


\begin{center}
 \begin{tabular}{rl}
  \vspace{10pt}
  &\xiaosan Name:八元\\
  \vspace{10pt}
  &\xiaosan No:\\
  \vspace{10pt}
  &\xiaosan Date:\\
  \vspace{10pt}
  &\xiaosan Score:\rule{3cm}{0.15mm}\\
  \end{tabular}
\end{center}
\vspace{100pt}

\begin{center}
\vspace{10pt}
{\xiaosan\textbf{SHANGHAI JIAOTONG UNIVERSITY}}\\   %宋体，小二，加粗
{\xiaosan\textbf{SJTU-ParisTech Elite Institute of Technology}} % Times New Roman，小二
\end{center}
% \maketitle

\newpage
\tableofcontents
\newpage


\section{Introduction}
\subsection{}
\subsection{}
\subsection{}

\section{Summary}

\bibliographystyle{plain}
\bibliography{ref}

\newpage

\begin{appendices}
\section{Code Example}
\lstinputlisting[language=Matlab]{code/ex.m}

\end{appendices}



\end{document}